% Datation par le carbone 14
\begin{minipage}{0.58\linewidth}
    \begin{Block}[Principe de la datation]
    \small
    Dans la haute atmosphère, les rayons cosmiques provoquent des réactions
    nucléaires qui libèrent des neutrons.

    Ces neutrons, une fois ralentis, sont absorbés par des noyaux d'azote
    \ch{_{7}^{14}N} au cours d'une réaction qui donne naissance à du carbone
    \ch{_{6}^{14}C} et à une autre particule.

    Le carbone 14 ainsi créé est radioactif.

    Le carbone 14 est alors assimilé comme le carbone stable (\ch{_{6}^{12}C})
    par les plantes au cours de la synthèse chlorophyllienne. Pendant toute leur
    vie, la proportion de carbone 14 reste très stable dans les plantes.

    À leur mort, la quantité de carbone 14 décroît par radioactivité. Il suffit
    donc de mesurer la proportion de carbone 14 restante dans l'échantillon
    étudié pour dater sa mort.
\end{Block}
\end{minipage}
\hfill
\begin{minipage}{0.42\linewidth}
En s'aidant éventuellement des documents ci-dessous, répondre aux questions
suivantes.

\vspace{3mm}

\begin{figure}[H]
    \centering
    \begin{tikzpicture}
        \def\No{100}
        \def\tm{5570}
        \pgfmathparse{\tm/ln(2)}\let\ta=\pgfmathresult
        \begin{axis}[
                cycle list name=my color list,
                width=8.5cm,height=7cm,
                xlabel={\small temps en années},
                ylabel={\small
                    \begin{tabular}{l}
                        nombre de noyaux de carbone radioactif \\
                        dans un échantillon de carbone
                    \end{tabular}},
                xmin=0,xmax=20000,
                ymin=0,ymax=120,
                xtick={0,\tm,2*\tm,3*\tm,4*\tm},
                ytick={0,\No/8,\No/4,\No/2,\No},
                variable=\t,
                samples=200,
                grid=none,
                scaled ticks=false,
                    xticklabel style={ /pgf/number format/.cd, 1000 sep={ } },
	            every axis x label/.append style={
                    at={([yshift=-3.5mm]current axis.right of origin)},
                    anchor=north east,inner xsep=0,
	            },
	            every axis y label/.append style={
                    at={(current axis.above origin)},
                    anchor=north west,inner ysep=0,
	            },
            ]
            \addplot[domain=0:19000,very thick] {\No*exp(-\t/\ta)};
        \end{axis}
    \end{tikzpicture}
\end{figure}
\end{minipage}


\begin{enumerate}
    \item Donner la composition des noyaux \ch{_{6}^{14}C} et \ch{_{6}^{12}C}.
          Comment appelle-t-on de tels noyaux~?
    \item Après avoir rappelé les lois de conservation, écrire l'équation de la
          réaction nucléaire relative au deuxième paragraphe du texte (réaction
          d'un neutron sur un atome d'azote).

          Quelle est la particule apparue en plus du carbone 14~?
    \item Le carbone 14 est radioactif $\beta^{-}$. Quelle est la nature de ce
          rayonnement? Écrire l'équation de la réaction nucléaire relative à la
          désintégration radioactive du carbone 14 .
    \item Qu'appelle-t-on demi-vie d'un noyau radioactif~?

          Quelle est la demi-vie du carbone 14~?
    \item Dans un échantillon de bois vivant, on détecte un atome de carbone 14
          pour $10^{12}$ atomes de carbone 12.
    
          Dans un morceau de bois ancien, mort, on constate qu'il n'y a plus
          qu'un atome de carbone 14 pour $8 \times 10^{12}$ atomes de carbone
          12. Quel est l'âge de ce morceau de bois~?
\end{enumerate}


